% Draft. Every number below is transcribed from docs/results-paper.md,
% docs/results-phase4.md and docs/results-phase7.md, which the scripts in scripts/paper/
% regenerate from committed runs. Each table names its script in a comment.
\documentclass[11pt]{article}

\usepackage[review]{acl}
\usepackage{times}
\usepackage{latexsym}
\usepackage[T1]{fontenc}
\usepackage[utf8]{inputenc}
\usepackage{microtype}
\usepackage{graphicx}
\usepackage{booktabs}
\usepackage{multirow}
\usepackage{amsmath}
\usepackage{pgfplots}
\pgfplotsset{compat=1.18}
\usepgfplotslibrary{groupplots}

\newcommand{\sys}{\textsc{graphwalk}}
\newcommand{\stopop}{\textsc{stop}}
\newcommand{\ci}[2]{{\scriptsize[#1,\,#2]}}

\title{When to Trust a Knowledge-Graph Walk:\\
An Empirical Study of Per-Hop Decision Confidence}

\author{Anonymous}

\begin{document}
\maketitle

\begin{abstract}
Question answering over knowledge graphs with large language models is accurate but
expensive, and it rarely says when an answer is likely wrong. We study an old design,
in which each hop of a walk is a classification over the current node's relations plus
\stopop{} \citep{zhang2022sr,das2018minerva}, with a zero-shot, swappable decider, and
ask what the walk's path confidence is worth. On MetaQA, WebQSP, CWQ and a
2WikiMultiHopQA evidence graph, we compare a hosted classifier, three open LLMs read
through their option-letter token probabilities, and LLM-stated confidence. Where walks
fail often enough to matter (MetaQA 3-hop, WebQSP), path confidence ranks right above
wrong answers well (AUROC 0.82--0.97): answering the most confident half of MetaQA
3-hop questions is 99\% accurate, against 81\% overall. With the model and prompt held
fixed, token probabilities beat stated per-option scores, stated top-1 confidence and
5-sample voting by 0.06--0.23 AUROC. The raw confidence is not a probability (ECE 0.26
on WebQSP), but a two-parameter Platt fit repairs much of that. Shuffling the options
changes an LLM decider's answer on 10--27\% of questions without moving aggregate
accuracy or AUROC. The walk is \emph{less} accurate than an LLM writing the relation
path in one call ($-0.08$ F1 on MetaQA 3-hop, $-0.15$ on WebQSP), cheaper only on
Freebase, and its confidence fails on 2Wiki and CWQ. As a tool for a strong agent it
cuts cost by 17\% at equal accuracy. We release the library, every run and the scripts
that produce every table.
\end{abstract}

% ======================================================================================
\section{Introduction}

LLM-based question answering over knowledge graphs (KGQA) now reaches high accuracy on
the standard Freebase benchmarks, by prompting a model to explore the graph
\citep{sun2024tog,jiang2025kgagent}, to generate relation paths \citep{luo2024rog}, or
to read retrieved subgraphs \citep{mavromatis2024gnnrag}. These systems make many LLM
calls per question, and none of them tells the caller, cheaply, which of its answers to
trust. A deployment that could answer the easy questions cheaply and confidently, and
pass the rest to an expensive system or a person, needs exactly that signal.

A KG walk offers one for free. If every hop is a choice among the relations leaving the
current node, plus \stopop{}, each choice comes with a probability distribution, and the
path's probability is a confidence for the answer it reaches. The design itself is old:
SR \citep{zhang2022sr} scores relations against an END relation and ranks paths by the
product of step probabilities, and MINERVA \citep{das2018minerva} walks with a policy
over outgoing edges. Both are trained per dataset. What has changed is that any LLM, or
a hosted classifier, can now make the per-hop choice zero-shot, which makes the decider
a component one can swap, and its confidence something one can compare across deciders.

This paper is an empirical study of that confidence, not a new method. We ask:
\begin{enumerate}
\setlength\itemsep{0pt}
\item Does path confidence separate right from wrong answers well enough for selective
  answering, and on which graphs? (\S\ref{sec:selective})
\item Is the signal a property of one decider, or of reading probabilities? With the
  model fixed, how does it compare with stated confidence and sampling consistency?
  (\S\ref{sec:signal})
\item Is it robust to the order the options are shown in, and is it calibrated?
  (\S\ref{sec:robust})
\item What does it buy in practice: against an LLM writing the query, in a cascade, and
  as an agent's tool? (\S\ref{sec:accuracy}, \S\ref{sec:use})
\end{enumerate}

Our findings are mixed, and we report them as such. Path confidence is informative
where walks fail often enough to need it (AUROC 0.82--0.97 on MetaQA 3-hop and WebQSP),
and the signal comes from reading token probabilities, not from a particular decider: an
open 27B model matches a dedicated classifier, and the same model's stated confidence is
clearly worse. But the walk is less accurate than one LLM call that writes the whole
relation path, it is cheaper only on Freebase's larger schemas, its raw confidence is
overconfident, and on two of five graphs (2Wiki, CWQ) the confidence does little. Our
contributions are:
\begin{itemize}
\setlength\itemsep{0pt}
\item To our knowledge, among the first evaluations of per-hop path probability as a
  selective-prediction signal for KGQA, across four decider types and five graphs, with
  the controls a reviewer of letter-reading LLMs would ask for: option order, format
  adherence, calibration, and a same-model stated-confidence baseline.
\item A measured cost and accuracy comparison against an LLM path writer and against
  multi-step RAG, including the settings where walking loses.
\item An agent-tool ablation with a memorization control: replacing entity names with
  aliases drops closed-book F1 from 0.42 to 0.01 and separates graph tools from text
  search by 0.25 F1, where with real names they tie.
\item An open-source library (\sys{}) and every run, so each table can be regenerated
  without new model calls.
\end{itemize}

% ======================================================================================
\section{Related Work}

\paragraph{Stepwise walkers and LLM traversal.} Relation-path walkers predate LLMs:
DeepPath \citep{xiong2017deeppath} and MINERVA \citep{das2018minerva} learn edge
policies, NSM \citep{he2021nsm} propagates over subgraphs, and SR \citep{zhang2022sr}
retrieves subgraphs by expanding relation paths with an END relation and ranks them by
path probability. Our walk is SR's in structure; the difference is that the decider is
zero-shot and replaceable. With LLMs, Think-on-Graph \citep{sun2024tog} prompts the
model at every hop (WebQSP/CWQ hits@1 76.2/58.9 with ChatGPT, 82.6/69.5 with GPT-4), RoG
\citep{luo2024rog} generates relation paths (85.7/62.6), StructGPT
\citep{jiang2023structgpt} and KG-Agent \citep{jiang2025kgagent} use tool interfaces,
and GNN-RAG \citep{mavromatis2024gnnrag} retrieves with a GNN; Pangu
\citep{gu2023pangu} argues for discriminating among candidates rather than generating.
Replacing the per-hop LLM with something cheaper is also known: EffiQA
\citep{dong2025effiqa}, KG-CoT \citep{zhao2024kgcot}, embedding-guided search
\citep{shrestha2025efficient} and RouterKGQA \citep{yuan2026routerkgqa}. We do not
compete with these systems on accuracy; our walker is greedy, single-entity and has no
constraint handling, and it scores far below them on WebQSP and CWQ.

\paragraph{Confidence in KGQA.} Ca2KG \citep{ren2026ca2kg} reports selective accuracy
against coverage on MetaQA and WebQSP with LLM-prompted confidence; CPR
\citep{lin2026cpr} learns path-level scores for conformal answer sets; UAG
\citep{ni2025uag} and DoublyCal \citep{lu2026doublycal} calibrate KG reasoning; SkewRoute
\citep{wang2025skewroute} routes between small and large LLMs on retrieval-score skew,
a special case of cascading \citep{dekoninck2025cascade}. Selective prediction and its
risk--coverage evaluation follow \citet{geifman2017selective} and, for QA,
\citet{kamath2020selective}. We differ in the signal (the decider's own per-hop
probabilities, no training) and in comparing it across deciders.

\paragraph{LLM confidence and multiple-choice bias.} Token probabilities are often
better calibrated than verbalized confidence \citep{kadavath2022know,xiong2024express,kapoor2024taught},
but \citet{tian2023just} find the opposite for RLHF-tuned API models judged by ECE, and
\citet{kim2026same} find the ranking flips with the protocol. Letter-choice LLMs prefer
some letters and positions regardless of content \citep{zheng2024mcq,pezeshkpour2024order},
and first-token probabilities can disagree with the text answer \citep{wang2024answerc};
\citet{zhuang2024beyond} read label probabilities for ranking. Our P6--P8 experiments
are designed around these results.

\paragraph{Memorization and agents.} WebQSP-style answers are partly memorized and
partly wrong \citep{zhang2025kgqagen}. ARoG \citep{ning2026arog} anonymizes WebQSP/CWQ
entities for privacy and finds that it breaks prompting methods; KG-LLM-Bench
\citep{markowitz2025kgllmbench} finds pseudonyms barely matter when the subgraph is in
context; \citet{longpre2021entity} substitute entities to study knowledge conflicts.
Tool-using agents over KGs \citep{gu2024middleware} and composite tools that save agent
calls \citep{abuzakuk2026awo} are the closest work to our agent ablation.

% ======================================================================================
\section{Method}
\label{sec:method}

\paragraph{The walk.} Given a question $q$ and a topic entity, the walk keeps a
\emph{frontier}, a set of nodes, starting from the topic entity. At each step the
options are the relations incident to the frontier, each with its direction, up to
five example target names, their types and a count, plus \stopop{} (``the current
nodes are what the query asks for; return them''). The decider returns a distribution
$p(\cdot \mid q, \text{path}, \text{options})$; the walk follows the chosen relation to
all of its targets (visited nodes excluded) and stops on \stopop{}, at a dead end, or at
depth 4 (three hops plus \stopop{}). The answer is the frontier at \stopop{}. All
results use the greedy setting (the argmax at each step); beam search is implemented
but not evaluated here. When the graph has node types, the decider also classifies the
expected answer type, which is shown at every step. Prompts are in
Appendix~\ref{app:prompts}.

\paragraph{Confidence.} A walk that made decisions with chosen-option probabilities
$p_1, \dots, p_L$ (the final \stopop{} included; forced moves with one legal option are
not decisions) gets the confidence
\begin{equation}
  c = \exp\Big(\frac{1}{L}\sum_{i=1}^{L} \log \max(p_i, 10^{-3})\Big),
\end{equation}
the geometric mean of the step probabilities. This is the length-normalized path
probability used to rank beams, not the raw product; the two order walks of equal
length identically. A walk with no answer gets $c = 0$.

\paragraph{Deciders.} We use four.
\textbf{Jev} is a hosted classification model (\citealp{jev}; version 1.13 through
OpenRouter) that takes a state, instructions and labeled options and returns a
probability per option; it is closed, and the main reason we test the others.
\textbf{Token probabilities} (\texttt{logprob}): a chat LLM sees the options lettered
A--T, answers with one token at temperature 0 with reasoning disabled, and the
distribution is read from that token's top log-probabilities, summed over each
letter's token variants and renormalized over the offered letters. Options beyond 20 are
ranked down by embedding similarity to the question (bge-small). We use Qwen3.8-27B
unless stated, and Gemma~4 31B and DeepSeek V4 Flash in one comparison.
\textbf{Stated scores} (\texttt{llm}): an LLM gives every option a 0--100 score, which
is normalized. In E3/P0 this is a different model (\texttt{gpt-6-luna}); in P8 the same
Qwen model and lettered prompt are used, and the reply format is the only change
(\S\ref{sec:signal}).

\paragraph{Escalation.} When $c$ falls below a threshold $t$, the question is passed to
a stronger, costlier system (the LLM decider, or the LLM path writer below), and the
escalated question pays for both.

% ======================================================================================
\section{Experimental Setup}
\label{sec:setup}

\paragraph{Data.} \textbf{MetaQA} \citep{zhang2018metaqa} 1-, 2- and 3-hop over the
WikiMovies graph (9 relations; answers are entity sets). \textbf{WebQSP}
\citep{yih2016webqsp} and \textbf{CWQ} \citep{talmor2018cwq} over the per-question
Freebase subgraphs released with RoG (median $\sim$4,400 triples), and, for WebQSP, one
graph merging all 1,628 test subgraphs (781k nodes, 2.28M edges, 5,419 relations).
\textbf{2Wiki}: the gold evidence triples of 2WikiMultiHopQA \citep{ho2020twowiki} dev
questions pooled into one graph, so that a question's own chain is not the whole graph.
Topic entities are given everywhere. Samples are seeded (seed 0): 200 questions per
MetaQA hop count (500 for P5), 300 for 2Wiki, 50--500 for WebQSP and 50 for CWQ, as
stated per table.

\paragraph{Baselines.} The \textbf{LLM path writer} sees the schema around the topic
entity (relations with endpoint types; for MetaQA, the 9 relations with one-line
glosses), writes a relation path in one call, which is executed like a walk, and may
retry twice when the path returns nothing. \textbf{Multi-step RAG} retrieves over a
textualized graph for up to 5 rounds. Both use \texttt{gpt-6-luna}.

\paragraph{Metrics.} Correctness is exact match of the answer set (EM) on MetaQA and
2Wiki and hits@1 on WebQSP/CWQ, as is standard there; we also report token-level F1 over
answer sets. Confidence is evaluated by AUROC (the chance that a right answer outranks a
wrong one), AURC (area under the risk--coverage curve; lower is better; ``random'' is
the error rate), expected calibration error (ECE, 10 equal-width bins) and Brier score.
Intervals are 95\% percentile bootstraps over questions (2,000 resamples); differences
between systems are paired. Cost is the provider-reported charge per 1,000 questions;
latency is wall time per question.

% ======================================================================================
\section{Accuracy and Cost Against an LLM Path Writer}
\label{sec:accuracy}

% Source: docs/results-paper.md (P5), docs/results-phase7.md (A2, A6),
% docs/results-phase4.md (E2b); scripts/paper/table_curated.py and scripts/eval/*.
\begin{table}[t]
\centering\small
\setlength\tabcolsep{3pt}
\resizebox{\columnwidth}{!}{%
\begin{tabular}{lrcccc}
\toprule
& & \multicolumn{2}{c}{F1} & \multicolumn{2}{c}{\$ / 1k q} \\
\cmidrule(lr){3-4}\cmidrule(lr){5-6}
Graph & $n$ & walk & path & walk & path \\
\midrule
MetaQA 2-hop & 500 & 0.985 & \textbf{0.999} & 0.096 & \textbf{0.061} \\
MetaQA 3-hop & 500 & 0.884 & \textbf{0.964} & 0.152 & \textbf{0.092} \\
WebQSP & 500 & 0.54 & \textbf{0.69} & \textbf{0.22} & 0.78 \\
WebQSP, one graph & 100 & 0.49 & \textbf{0.66} & \textbf{0.24} & 1.48 \\
CWQ & 50 & 0.32 & 0.29 & \textbf{0.22} & 1.23 \\
\midrule
2Wiki, text-extracted & 120 & \textbf{0.254} & 0.079 & \textbf{0.18} & 3.63 \\
HotpotQA, text-extracted & 120 & \textbf{0.182} & 0.040 & \textbf{0.17} & 6.09 \\
\bottomrule
\end{tabular}}
\caption{Jev walk vs.\ an LLM writing the relation path (2 retries). On curated graphs
the path writer is more accurate; the walk is cheaper only on Freebase. On graphs
extracted from text by an LLM (bottom), the written path rarely matches an existing
relation, but every graph method there is far below multi-step RAG over the text
(0.75--0.82 F1).}
\label{tab:accuracy}
\end{table}

Table~\ref{tab:accuracy} sets the stakes for the rest of the paper: as a question
answerer, the walk is not the better system. On MetaQA's 9-relation schema, one LLM call
that writes the path is more accurate (+0.08 F1 at 3 hops) and \emph{cheaper}; the walk
is only faster (p50 1.1\,s vs.\ 2.2\,s). On WebQSP the path writer is 15 F1 points
better and the walk 3.6$\times$ cheaper, with a much shorter tail (p95 3.1\,s vs.\
15.8\,s). A schema too large to read does not rescue walking on a curated graph: on the
merged 5,419-relation WebQSP graph, giving the path writer only the relations within two
hops of the topic entity keeps it 16 F1 points ahead (+0.16 \ci{+0.07}{+0.25}). On CWQ
every system is weak (0.29--0.33 F1, against hits@1 of 0.59--0.72 reported by
recent LLM methods), because
none handles CWQ's constraints and a fifth of the sampled questions have no answer in
their subgraph; parity there shows nothing. Walking wins only on noisy graphs extracted
from text, where an LLM's plausible relation names usually do not exist (81\% of its
2Wiki paths returned nothing), but there multi-step RAG over the source text is far
better than either.

The case for the walk is therefore not accuracy. It is that each answer arrives with a
confidence, cheaply. The rest of the paper asks how good that confidence is.

% ======================================================================================
\section{Path Confidence as a Selective-Prediction Signal}
\label{sec:selective}

% Source: scripts/paper/figure_selective.py (P0/P1), docs/results-paper.md.
\begin{table*}[t]
\centering\small
\begin{tabular}{lr ccc ccc ccc}
\toprule
& & \multicolumn{3}{c}{Accuracy} & \multicolumn{3}{c}{AUROC} & \multicolumn{3}{c}{AURC (random order)} \\
\cmidrule(lr){3-5}\cmidrule(lr){6-8}\cmidrule(lr){9-11}
Graph & $n$ & Jev & stated & token & Jev & stated & token & Jev & stated & token \\
\midrule
MetaQA 1-hop & 200 & 0.93 & 0.95 & 0.94 & 0.64 & 0.66 & 0.58 & .040 (.070) & .030 (.055) & .043 (.065) \\
MetaQA 2-hop & 200 & 0.98 & 0.99 & 0.97 & \textbf{0.94} & 0.54 & \textbf{0.97} & .002 (.020) & .008 (.015) & .001 (.030) \\
MetaQA 3-hop & 200 & 0.81 & 0.82 & 0.80 & \textbf{0.97} & 0.64 & \textbf{0.96} & .028 (.195) & .149 (.180) & .032 (.200) \\
2Wiki evidence graph & 300 & 0.88 & 0.93 & 0.85 & 0.67 & 0.69 & \textbf{0.77} & .141 (.117) & .037 (.073) & .139 (.147) \\
WebQSP & 50 & 0.54 & 0.66 & 0.62 & \textbf{0.87} & 0.64 & \textbf{0.84} & .189 (.460) & .238 (.340) & .154 (.380) \\
CWQ & 50 & 0.28 & 0.28 & 0.28 & 0.59 & 0.54 & \textbf{0.73} & .611 (.720) & .750 (.720) & .607 (.720) \\
\bottomrule
\end{tabular}
\caption{Path confidence as a selective-prediction signal, for the same walks with three
deciders: Jev, an LLM's stated 0--100 scores (\texttt{gpt-6-luna}), and an open LLM's
token probabilities (Qwen3.8-27B). Accuracy is EM on MetaQA and 2Wiki, hits@1 on
WebQSP/CWQ. AURC: lower is better; in parentheses, the AURC of a random order (the error
rate). Bold: AUROC $\geq 0.70$.}
\label{tab:selective}
\end{table*}

Table~\ref{tab:selective} compares three deciders on the same walks. Three patterns
stand out.

\paragraph{Where walks fail often enough to matter, probability-based confidence is
strongly informative.} On MetaQA 3-hop (about one walk in five is wrong), Jev and the
open model's token probabilities reach AUROC 0.97 and 0.96, and AURC 0.028--0.032
against 0.20 for a random order: answering the most confident half is 99\% accurate,
against 81\% overall. On WebQSP, the confidence stays informative at ten times the
pilot's sample (Jev, $n=500$: AUROC 0.86 with EM). On MetaQA 1--2 hop, 93--98\% of walks
are right and there is little to rank; AUROC is noisy there (0.54--0.97, from 2--14 wrong
answers).

\paragraph{The signal is not Jev's.} The open model's letter probabilities match Jev
on MetaQA 2--3 hop and WebQSP and are better on 2Wiki and CWQ, at similar accuracy. Jev
is about $10\times$ faster per walk (p50 1.1\,s vs.\ 9.9\,s on MetaQA 3-hop through
OpenRouter) and cheaper (\$0.15 vs.\ \$0.26 per 1k questions). The model does matter: on
MetaQA 3-hop, Gemma~4 31B is as accurate (0.82 EM) but overconfident (mean confidence
0.995; AUROC 0.82), and DeepSeek V4 Flash rarely chooses \stopop{} and collapses to 0.44
EM, although its confidence still ranks (AUROC 0.84; Appendix~\ref{app:models}).

\paragraph{Stated scores are weak, and the confidence fails on some graphs.} The
stated-score LLM decider is as accurate or more, but its confidence barely ranks
answers (AUROC 0.54--0.69). On 2Wiki, sorting by Jev's confidence is \emph{worse} than a
random order (AURC 0.141 vs.\ 0.117) despite AUROC 0.67: a few confidently wrong walks
sit at the top. On CWQ ($n=50$, 28\% correct) the best AUROC is 0.73 and no decider's
confidence is useful after recalibration (\S\ref{sec:robust}). Path confidence measures
whether the decider was sure of each relation choice; it cannot see a constraint the
walk ignored, or a graph that lacks the fact.

% ======================================================================================
\section{Is It the Signal or the Model?}
\label{sec:signal}

% Source: scripts/paper/table_stated.py (P8); deltas: paired bootstrap, docs/results-paper.md.
\begin{table}[t]
\centering\small
\setlength\tabcolsep{2.5pt}
\resizebox{\columnwidth}{!}{%
\begin{tabular}{lcccc}
\toprule
Confidence signal & Acc. & AUROC & $\Delta$ AUROC & \$/1k \\
\midrule
\multicolumn{5}{l}{\emph{MetaQA 3-hop (EM), $n=200$}} \\
token probabilities & 0.800 & \textbf{0.955} & -- & 0.26 \\
stated, per option & 0.740 & 0.719 & $-0.23$ \ci{$-0.33$}{$-0.14$} & 0.63 \\
stated, top-1 & 0.795 & 0.748 & $-0.21$ \ci{$-0.30$}{$-0.12$} & 0.38 \\
vote, 5 samples & 0.785 & 0.894 & $-0.06$ \ci{$-0.11$}{$-0.02$} & 1.49 \\
\midrule
\multicolumn{5}{l}{\emph{WebQSP (hits@1), $n=200$}} \\
token probabilities & 0.615 & \textbf{0.821} & -- & 0.29 \\
stated, per option & 0.560 & 0.736 & $-0.09$ \ci{$-0.17$}{$-0.01$} & 0.76 \\
stated, top-1 & 0.600 & 0.695 & $-0.13$ \ci{$-0.20$}{$-0.06$} & 0.40 \\
vote, 5 samples & 0.565 & 0.767 & $-0.06$ \ci{$-0.12$}{$+0.01$} & 1.67 \\
\bottomrule
\end{tabular}}
\caption{Same model (Qwen3.8-27B), endpoint and lettered prompt; only the confidence
signal changes. $\Delta$: paired bootstrap against token probabilities. A plain rerun of
the token-probability walk moves AUROC by $+0.009$ \ci{$-0.005$}{$+0.025$}.}
\label{tab:stated}
\end{table}

\begin{figure}[t]
\centering
\begin{tikzpicture}
\begin{groupplot}[
  group style={group size=2 by 1, horizontal sep=0.9cm, y descriptions at=edge left},
  width=4.6cm, height=4.4cm, xmin=0, xmax=1, ymin=0,
  xlabel={coverage}, tick label style={font=\scriptsize}, label style={font=\scriptsize},
  title style={font=\scriptsize, yshift=-1ex}, no markers, every axis plot/.append style={thick},
  legend style={font=\tiny, draw=none, at={(1.05,-0.42)}, anchor=north, legend columns=4},
  legend cell align=left,
]
\nextgroupplot[title={MetaQA 3-hop}, ylabel={risk (1 $-$ accuracy)}, ymax=0.3]
\addplot[black] table[x=coverage, y=risk] {figures/selective-metaqa3hop-token.dat};
\addplot[blue, dashed] table[x=coverage, y=risk] {figures/selective-metaqa3hop-scores.dat};
\addplot[red, dotted] table[x=coverage, y=risk] {figures/selective-metaqa3hop-top1.dat};
\addplot[teal, dashdotted] table[x=coverage, y=risk] {figures/selective-metaqa3hop-vote.dat};
\legend{token probs, stated/option, stated top-1, vote}
\nextgroupplot[title={WebQSP}]
\addplot[black] table[x=coverage, y=risk] {figures/selective-webqsp-token.dat};
\addplot[blue, dashed] table[x=coverage, y=risk] {figures/selective-webqsp-scores.dat};
\addplot[red, dotted] table[x=coverage, y=risk] {figures/selective-webqsp-top1.dat};
\addplot[teal, dashdotted] table[x=coverage, y=risk] {figures/selective-webqsp-vote.dat};
\end{groupplot}
\end{tikzpicture}
\caption{Risk--coverage curves for the signals of Table~\ref{tab:stated}: the walk
answers only its most confident questions, down to each coverage. Lower is better.}
\label{fig:riskcov}
\end{figure}

The comparison in Table~\ref{tab:selective} confounds two things: the stated scores came
from a different model. We therefore fix the model (Qwen3.8-27B), the endpoint and the
lettered prompt, and change only the confidence signal: token probabilities; a 0--100
score per option; one letter with a 0--100 confidence in it, the verbalized top-1
format of \citet{tian2023just} (the remaining mass spread evenly); and the vote share of
5 samples at temperature 1 (sampling consistency, with 0.1 pseudo-votes per option).
Table~\ref{tab:stated} and Figure~\ref{fig:riskcov} give the result.

\paragraph{With the model fixed, the signal is what matters.} Both stated formats lose
0.09--0.23 AUROC to token probabilities, every interval excluding zero. Asking for one
confidence in the chosen answer, the stronger stated baseline, is no better than scores
per option. Sampling consistency is the best alternative but still worse ($-0.06$; on
WebQSP the interval touches zero), it costs 5--6$\times$ as much per decision, and
sampling at temperature 1 costs some accuracy. Token probabilities are also the cheapest
LLM signal: one output token per decision, 1.4--2.6$\times$ cheaper per question than
the stated formats. Parsing was not the cause: 1.6\% (MetaQA) and 3.4\% (WebQSP) of
per-option replies could not be parsed and were scored uniformly; the other formats
parsed in all but one case.

This agrees with \citet{kadavath2022know} and \citet{xiong2024express}, and differs from
\citet{tian2023just}, who found verbalized confidence better calibrated for RLHF-tuned
API models judged by ECE. We tested one open model on two datasets and judge mainly by
ranking; we cannot test their API models, most of which do not expose log-probabilities.

% ======================================================================================
\section{Robustness and Calibration}
\label{sec:robust}

\paragraph{Option order.} The walk shows options in a fixed order (by direction and
relation name, \stopop{} last), and letter-reading LLMs are known to prefer letters and
positions \citep{zheng2024mcq,pezeshkpour2024order}. We reran MetaQA 3-hop and WebQSP
($n=200$) with each question's options in three random orders, next to a plain rerun in
the original order to measure noise at temperature 0 (Appendix~\ref{app:order}).
\emph{In aggregate, order barely matters}: accuracy moves within 2.5 points and AUROC
stays at 0.89--0.95 (MetaQA) and 0.80--0.82 (WebQSP). Of six paired AUROC changes for
Qwen, one is borderline ($-0.064$ \ci{$-0.132$}{$-0.000$}) and the rest include zero.
\emph{Per question, it matters}: a new order changes Qwen's answer on 10--16\% of MetaQA
and 22--27\% of WebQSP questions, against 4.5\% and 10\% for a plain rerun, and the rank
correlation of confidence across orders falls from 0.90 to 0.61--0.73 on MetaQA. Jev is
far more stable (6\% vs.\ 1.5\%). The confidence says which answers are right on
average; it is not a stable property of each question. Agreement across orders is a
worse confidence than the path probability (AUROC 0.91 vs.\ 0.955 on MetaQA, 0.76 vs.\
0.82 on WebQSP), so averaging over orders is not an obvious fix.

\paragraph{Format adherence.} 96--97\% of the first token's probability falls on
offered letters, the top token was not a letter in at most 1.2\% of decisions, and no
decision lacked a letter entirely, so renormalizing over letters discards little.

% Source: scripts/paper/table_recalibration.py (P7).
\begin{table}[t]
\centering\small
\setlength\tabcolsep{3pt}
\resizebox{\columnwidth}{!}{%
\begin{tabular}{llcccc}
\toprule
Graph & Decider & Acc. & Conf. & ECE raw$\to$Platt & Skill \\
\midrule
\multirow{3}{*}{MetaQA 3-hop} & Jev & 0.81 & 0.92 & 0.14 $\to$ 0.08 & \textbf{+0.62} \\
& stated & 0.82 & 0.97 & 0.15 $\to$ 0.07 & +0.02 \\
& token & 0.80 & 0.93 & 0.13 $\to$ 0.07 & \textbf{+0.57} \\
\midrule
\multirow{3}{*}{2Wiki} & Jev & 0.88 & 0.94 & 0.05 $\to$ 0.05 & +0.16 \\
& stated & 0.93 & 0.98 & 0.06 $\to$ 0.03 & +0.33 \\
& token & 0.85 & 0.96 & 0.11 $\to$ 0.07 & +0.35 \\
\midrule
\multirow{2}{*}{WebQSP} & Jev (500) & 0.54 & 0.80 & 0.26 $\to$ 0.10 & +0.28 \\
& token (200) & 0.62 & 0.84 & 0.22 $\to$ 0.12 & +0.25 \\
\midrule
CWQ (50) & all three & 0.28 & 0.67--0.72 & 0.39--0.44 $\to$ 0.18--0.21 & $\leq -0.02$ \\
\bottomrule
\end{tabular}}
\caption{Calibration before and after Platt scaling (logistic fit of correctness on
$\log c$), cross-fitted on halves, 10 splits. Conf.: mean raw confidence. Skill:
$1 - \text{Brier}/\text{Brier}_{\text{base rate}}$ after scaling (0 = no better than
knowing the accuracy).}
\label{tab:calibration}
\end{table}

\paragraph{Calibration.} AUROC ignores scale, and \citet{tian2023just} and
\citet{kim2026same} judge confidence by calibration error instead.
Table~\ref{tab:calibration} reports ECE and Brier skill before and after Platt scaling
\citep{platt1999,guo2017calibration}, fit on half the questions and scored on the other
half. \emph{Raw confidence is not a probability}: wherever accuracy is below about 90\%
it is overconfident, badly on Freebase (mean 0.80 for 54\% correct on WebQSP). \emph{A
two-parameter fit on a few hundred labeled answers repairs much of it} (WebQSP ECE 0.26
$\to$ 0.10, skill +0.25--0.28). After scaling, the decider ordering of
\S\ref{sec:selective} holds for calibration too: on MetaQA 3-hop, Jev and token
probabilities reach skill +0.62 and +0.57 and the stated scores +0.02, although their
raw ECE is similar (0.13--0.15). Raw ECE alone would have hidden the difference. Where
the confidence failed to rank (2Wiki, CWQ), scaling cannot rescue it. ECE with 10 bins
is noisy at these sample sizes, and each fit sees only half the questions.

% ======================================================================================
\section{What the Confidence Buys}
\label{sec:use}

\paragraph{Escalation.} Simulated offline from runs in which both systems answered the
same questions (Appendix~\ref{app:escalation}): on MetaQA 3-hop, passing walks with
$c < 0.9$ (25\% of questions) to the LLM decider recovers all of its +0.06 F1 lead at
\$0.30 instead of \$0.55 per 1k questions. On 2Wiki, where the confidence is weaker,
closing the gap takes $t=0.95$ (34\% escalated, \$0.13 vs.\ \$0.16). On Freebase,
escalation saves little, because the questions the walk answers confidently are the easy
ones that the path writer also answers: on WebQSP ($n=50$), matching the path writer
means escalating two thirds of the questions at 96\% of its cost, and on the merged
graph a cascade reaches its accuracy for 19\% less. Escalation pays on small curated
schemas, where the expensive system is a per-hop LLM, not where a single LLM call writes
the path.

\paragraph{Can one tell before walking?} A logistic regression on pre-walk features
(question length and cues, entity degree, relation counts) predicts walk success with
AUROC 0.41--0.68 within a graph, against 0.70--0.96 for the walk's confidence; adding the
features to the confidence does not help. A walk costs about \$0.0002 and a second, so
the routing rule is to walk first and read the confidence.

% Source: scripts/paper/table_agent.py (P4), scripts/paper/table_anon.py (P3).
\begin{table}[t]
\centering\small
\setlength\tabcolsep{3pt}
\resizebox{\columnwidth}{!}{%
\begin{tabular}{lcc}
\toprule
\multicolumn{3}{l}{\emph{(a) Strong agent (Gemini 3.1 Pro), WebQSP, $n=100$}} \\
Arm & F1 & \$/q \\
\midrule
graph tools & 0.747 \ci{0.672}{0.823} & 0.0376 \\
graph tools + walk & 0.755 \ci{0.683}{0.824} & 0.0311 \\
$\Delta$ & $+0.008$ \ci{$-0.038$}{$+0.057$} & $-17\%$ \ci{$-27\%$}{$-8\%$} \\
\midrule
\multicolumn{3}{l}{\emph{(b) Cheap agent, F1 with real names $\to$ aliases, $n=100$}} \\
Arm & real & aliases \\
\midrule
closed book & 0.42 & 0.01 \\
walk alone (no agent) & 0.51 & 0.36 \\
agent + graph tools & 0.74 & 0.45 \\
agent + graph tools + walk & 0.72 & \textbf{0.52} \\
agent + text search & 0.68 & 0.20 \\
\bottomrule
\end{tabular}}
\caption{The walk as an agent's tool. (a) Paired difference of the walk arm. (b) Entity
names replaced by meaningless aliases in graph, gold answers and questions. With aliases,
graph tools beat search by $+0.25$ \ci{$+0.16$}{$+0.35$} F1, graph tools + walk by
$+0.32$ \ci{$+0.23$}{$+0.42$}.}
\label{tab:agent}
\end{table}

\paragraph{As an agent's tool.} One tool-calling loop, the same prompt and a 10-turn
budget for every arm; only the tools differ: graph exploration (\texttt{relations},
\texttt{neighbors}), the same plus \texttt{walk} (which returns answers with their
confidence), or dense search over the same facts as text. With a strong agent
(Table~\ref{tab:agent}a), adding \texttt{walk} cuts cost by 17\% at equal F1. The
saving is smaller than in our first 30 questions ($-25\%$), and its mechanism is not the
one we expected: the agent does not stop at a confident walk but verifies it, spending
fewer exploration calls and, in the first 30 questions, about 40\% less reasoning
output. With a cheap agent, the
walk's own decision cost cancels the token savings.

\paragraph{A memorization control.} The cheap agent answers 42\% of WebQSP F1 closed
book, so any tool comparison partly measures memory. We replaced every entity name with
a stable alias (``Entity 3fa9c1'') in the graph, the gold answers and the question text,
keeping relations, numbers and dates (Table~\ref{tab:agent}b). Closed book falls to
0.01, so the control works. Every arm leaned on names, text search the most ($-0.48$
F1): without them, an agent exploring the graph beats one searching the same facts by
0.25--0.32 F1, while with real names they were within 0.06. Part of this is dense
retrieval losing meaningful names to embed, so the aliased graph is a lower bound for
search rather than a model of a private KG. The walk alone loses 0.15 F1, because its
decider also read what the names mean. Unlike KG-LLM-Bench
\citep{markowitz2025kgllmbench}, which finds pseudonyms barely matter with the subgraph
in context, the questions here are ones the model may have memorized.

% ======================================================================================
\section{Negative Results}
\label{sec:negative}

We report the settings where walking did not help, since each was a plausible use.
\textbf{QA over text}: on 2Wiki and HotpotQA paragraphs ingested into a graph by LLM
extraction, walking loses to multi-step RAG over the same text by 10--19 F1 points, and
a 20$\times$ pricier extraction model does not close the gap. \textbf{As a RAG add-on},
the graph's first retrieval cut multi-step RAG's LLM rounds by 16--21\% but gave no
reliable accuracy gain, and the classifier calls cost more than the saved LLM calls.
\textbf{Large curated schemas} do not favor walking (\S\ref{sec:accuracy}).
\textbf{Constraints}: CWQ needs filters and superlatives that a relation walk cannot
express. \textbf{Pre-walk routing}: question features do not predict success
(\S\ref{sec:use}).

% ======================================================================================
\section{Conclusion}

Reading a decider's probabilities at each hop of a KG walk gives a confidence that,
where walks fail often enough to need one, separates right from wrong answers well; it
does not depend on a particular classifier; and, with the model fixed, it beats asking
the model how sure it is. It is not a probability until recalibrated, it is unstable per
question under option reordering, and it fails on graphs where errors come from missing
facts or ignored constraints. The walk it comes from is cheaper but less accurate than
an LLM writing the query. We see its use as a cheap first pass whose confidence decides
what to escalate, and as a tool that makes strong agents cheaper, not as a competitor
to LLM-based KGQA on accuracy.

\section*{Limitations}

\textbf{Samples.} 200--500 questions per curated setting, but 50 on the WebQSP and CWQ
pilots of Table~\ref{tab:selective} and 100 for the agent experiments; one seed for
most runs (three for the MetaQA headline comparisons). \textbf{Models.} Hosted models
whose versions can change behind the same name; Jev is closed. The open-model results
come from one provider's serving of each model. \textbf{Scope.} Topic entities are
given; the walker is greedy over relations from one entity, with no beam search,
constraint handling or multi-entity questions, which caps accuracy on WebQSP and CWQ.
Freebase is the only large curated graph. \textbf{Calibration protocol.} We judge
confidence mostly by ranking; Platt scaling needs labeled answers from the target graph.
\textbf{Data quality.} A large share of WebQSP/CWQ-style gold answers may be wrong or
incomplete \citep{zhang2025kgqagen}, which bounds every accuracy number.

\section*{Ethics and Disclosure}

The code, the experiments and a first draft of this paper were produced with an AI
assistant; the author reviewed the results and the text and is responsible for the
claims. The datasets are public. The
memorization control replaces entity names with aliases and introduces no personal data.

\bibliography{refs}

% ======================================================================================
\appendix

\section{Prompts}
\label{app:prompts}

\paragraph{Token-probability decider.} System: ``You answer one multiple-choice
question about the state given. Read the instructions, then reply with the letter of the
single best option and nothing else.'' User message:
\begin{small}
\begin{verbatim}
State:
{"query": "<question>"}

Instructions:
{"task": "You are walking a knowledge graph to
  answer the query in the state. If the current
  nodes are the things the query asks for, choose
  STOP. Otherwise choose the relation to follow
  from the current nodes that moves one step
  closer to the answer.",
 "current": {"nodes": [...], "count": N,
             "types": [...]},
 "path_so_far": [...],
 "the_query_asks_for_a": "<type>"}

Options:
A. o1: {"relation": "directed_by",
        "direction": "incoming",
        "leads_to": [...], "count": N,
        "leads_to_types": [...]}
...
F. STOP: "The current nodes are what the
          query asks for. Return them as the
          answer."

Answer with one letter.
\end{verbatim}
\end{small}
Jev receives the same state, instructions and options as a labeled choice question.
The stated-confidence variants (P8) replace only the last line, with ``Score every
option from 0 (certainly wrong) to 100 (certainly right), one per line as LETTER=SCORE
(for example A=80), and nothing else.'' or ``Reply with the letter of the single best
option, a space, and how confident you are that it is right, from 0 to 100 (for example:
B 85). Nothing else.'', and their system prompt asks for ``exactly the format asked
for''.

\section{More Open Models}
\label{app:models}

% Source: docs/results-paper.md (P1b).
\begin{table}[h]
\centering\small
\setlength\tabcolsep{3pt}
\resizebox{\columnwidth}{!}{%
\begin{tabular}{lrcccc}
\toprule
Decider & $n$ & EM & AUROC & ECE & \$/1k \\
\midrule
Jev & 500 & 0.80 & 0.95 & 0.11 & 0.15 \\
Qwen3.8-27B & 200 & 0.80 & 0.96 & 0.13 & 0.26 \\
Gemma 4 31B & 200 & 0.82 & 0.82 & 0.18 & 0.30 \\
DeepSeek V4 Flash & 200 & 0.44 & 0.84 & 0.36 & 0.09 \\
\bottomrule
\end{tabular}}
\caption{MetaQA 3-hop with three open models as token-probability deciders.}
\end{table}

\section{Option Order}
\label{app:order}

% Source: scripts/paper/table_order.py (P6).
\begin{table}[h]
\centering\small
\setlength\tabcolsep{2.5pt}
\resizebox{\columnwidth}{!}{%
\begin{tabular}{llccc}
\toprule
Setting & Run & Acc. & AUROC & Same ans. \\
\midrule
MetaQA 3-hop & original & 0.800 & 0.955 & -- \\
\quad Qwen & rerun & 0.795 & 0.963 & 0.955 \\
& 3 shuffles & .775--.810 & .890--.947 & .840--.895 \\
MetaQA 3-hop & original & 0.805 & 0.965 & -- \\
\quad Jev & rerun & 0.805 & 0.968 & 0.985 \\
& 1 shuffle & 0.805 & 0.955 & 0.940 \\
WebQSP & original & 0.615 & 0.821 & -- \\
\quad Qwen & earlier run & 0.620 & 0.839 & 0.900 \\
& 3 shuffles & .600--.625 & .798--.817 & .735--.780 \\
\bottomrule
\end{tabular}}
\caption{Accuracy and AUROC under shuffled option orders; ``same ans.'': share of
questions with the same answer set as the original run.}
\end{table}

\section{Escalation}
\label{app:escalation}

% Source: scripts/paper/table_escalation.py (A1), docs/results-phase7.md (A2, A6).
\begin{table}[h]
\centering\small
\setlength\tabcolsep{3pt}
\resizebox{\columnwidth}{!}{%
\begin{tabular}{llccc}
\toprule
Graph & System & Esc. & F1 & \$/1k \\
\midrule
MetaQA 3-hop & Jev only & 0\% & 0.897 & 0.154 \\
& LLM decider only & 100\% & 0.956 & 0.548 \\
& escalate $c<0.9$ & 25\% & 0.956 & 0.303 \\
2Wiki & Jev only & 0\% & 0.915 & 0.048 \\
& LLM decider only & 100\% & 0.941 & 0.164 \\
& escalate $c<0.95$ & 34\% & 0.942 & 0.126 \\
WebQSP (50) & path writer only & 100\% & 0.76 & 0.72 \\
& Jev $\to$ path, $c<0.95$ & 66\% & 0.76 & 0.70 \\
WebQSP, one graph & path writer only & 100\% & 0.66 & 1.48 \\
& Jev $\to$ path, $c<0.95$ & 68\% & 0.66 & 1.20 \\
\bottomrule
\end{tabular}}
\caption{Escalation simulated offline; an escalated question pays for both systems.}
\end{table}

\section{Reproducibility}

Every run is committed with its per-question scores, and each table is printed by a
script in \texttt{scripts/paper/} from those scores, without model calls. Re-running an
experiment needs an OpenRouter key and costs the amounts stated in the repository's
results pages (under \$2 per experiment except the strong-agent runs, about \$8 in
total).

\end{document}
